\chapter{Test curves and Legendre duality}
\label{chap:hermitian-testcurves}

A concave family of singularity envelopes determines a ray by taking its Legendre transform. The inverse transform extracts singular potentials from the ray, even when every finite slice of the ray is bounded. The correspondence rests on two maximality properties: an envelope fixed point cannot be increased within its singularity type below zero, and a Perron segment cannot be increased while keeping its endpoints. Legendre duality turns a finite-interval replacement of the ray into a bounded change of each singular potential.

Fix $\varphi\in\PSH(X,\theta)$ with $P_\theta[\varphi]=\varphi$. Thus $\varphi\leq0$ and $\sup_X\varphi=0$. We use the pointwise transforms of \cref{sec:psh-legendre} and the Perron segments of \cref{sec:perron-segments}. The envelope argument follows the treatment of the Ross--Witt Nystr\"om correspondence in \cite[chapter on test curves]{SGPT}.

\section{Envelope test curves and rays}

The cutoff of a test curve will become the linear growth rate of its ray. Normalizing each member by the singularity envelope makes this relation exact.

\begin{definition}\label{def:hermitian-envelope-testcurve}
An \emph{envelope test curve} with initial potential $\varphi$ consists of a real number $c$ and potentials $\Gamma_\tau\in\PSH(X,\theta)$ for $\tau<c$ such that
\begin{enumerate}
\item $\tau\mapsto\Gamma_\tau$ is pointwise concave and decreasing;
\item $P_\theta[\Gamma_\tau]=\Gamma_\tau$ for every $\tau<c$;
\item $(\sup_{\tau<c}\Gamma_\tau)^*=\varphi$.
\end{enumerate}
Its cutoff is $c$. We extend it by
\[
 \Gamma_c=\lim_{\tau\uparrow c}\Gamma_\tau,
 \qquad \Gamma_\tau\equiv-\infty\quad(\tau>c).
\]
\end{definition}

Each $\Gamma_\tau$ has supremum zero. Consequently \cref{thm:normalized-compactness,cor:decreasing-limit} show that $\Gamma_c$ belongs to $\PSH(X,\theta)$. The regularized supremum at the other end of the parameter range specifies the initial potential.

\begin{definition}\label{def:hermitian-envelope-ray}
An \emph{envelope ray} with initial potential $\varphi$ and slope $c$ is a subgeodesic $u_t$, $t>0$, such that
\begin{enumerate}
\item $u_t\to\varphi$ in $L^1(X)$ as $t\downarrow0$, and every finite restriction, with $u_0=\varphi$, is a Perron segment;
\item $P_\theta[u_t]=\varphi$ for every $t>0$;
\item $\sup_Xu_t=ct$ for every $t>0$.
\end{enumerate}
\end{definition}

The forward and inverse operations are \eqref{eq:hermitian-legendre-forward} and \eqref{eq:hermitian-legendre-backward}. By \cref{thm:hermitian-legendre-usc}, the forward transform of a test curve is a subgeodesic with its pointwise values unchanged. The remaining conditions on an envelope ray express the effect of the envelope fixed points.

\section{From a test curve to a ray}

The normalization of a test curve gives the supremum and initial value of its transform. Perron maximality takes more work: replacing the transform on a finite interval changes each inverse transform by a bounded amount, so the envelope fixed-point condition forces the replacement to agree with the original ray.

\begin{theorem}\label{thm:testcurve-to-envelope-ray}
The forward transform of an envelope test curve is an envelope ray with the same initial potential and cutoff slope. Moreover $U_t-ct$ decreases as $t$ increases.
\end{theorem}
\begin{proof}
By \cref{thm:hermitian-legendre-usc}, $U_t$ is a subgeodesic. The inequalities
\[
 \Gamma_\tau+t\tau\leq U_t\leq\varphi+ct
\]
give $\sup_XU_t\geq t\tau$ for each $\tau<c$ and $\sup_XU_t\leq ct$. Letting $\tau\uparrow c$ proves equality. The formula
\[
 U_t-ct=\sup_{\tau<c}\bigl(\Gamma_\tau-t(c-\tau)\bigr)
\]
proves decreasingness.

The order property of the singularity envelope gives
\[
 \Gamma_\tau=P_\theta[\Gamma_\tau]\leq P_\theta[U_t]
 \leq P_\theta[\varphi]=\varphi.
\]
Taking the regularized supremum in $\tau$ proves $P_\theta[U_t]=\varphi$.

The supremum bounds and \cref{thm:normalized-compactness} give $L^1$ compactness as $t\downarrow0$. Any cluster point $v$ satisfies $v\leq\varphi$ by the upper bound and $v\geq\Gamma_\tau$ for every fixed $\tau<c$ by the lower bound. Here the inequalities pass first almost everywhere along a convergent subsequence, then everywhere by \cref{lem:qpsh-almost-everywhere-order}. Taking the regularized supremum in $\tau$ yields $v=\varphi$. Thus $U_t\to\varphi$ in $L^1$.

It remains to prove Perron maximality. Fix $0\leq a<b$ and replace $U_t$ on $(a,b)$ by the Perron segment between $U_a,U_b$. It exists and attains its endpoints by \cref{prop:hermitian-perron}, since the restricted original path is a competitor. By \cref{lem:qpsh-gluing,prop:subgeodesic-endpoint-control}, the resulting path $v_t$ is a subgeodesic, equals $U_t$ outside $(a,b)$, and dominates it inside. Convexity on the replaced interval gives $v_t\leq ct$ everywhere.

Fix $\tau<c$ and write
$\chi_\tau=\inf_{t>0}(v_t-t\tau)$. By \cref{prop:hermitian-kiselman} it is $\theta$-plurisubharmonic: the lower bound by $\Gamma_\tau$ excludes collapse to $-\infty$. We have
\begin{equation}\label{eq:finite-replacement-dual-bound}
 \Gamma_\tau\leq\chi_\tau
 \leq\Gamma_\tau+(b-a)(c-\tau).
\end{equation}
Indeed for $s>0$ with $s\in[a,b]$, the decreasingness of $U_t-ct$ gives
\[
 \chi_\tau\leq U_b-b\tau
 \leq U_s-s\tau+(b-s)(c-\tau).
\]
For $s$ outside the interval, use $v_s=U_s$ to obtain the same bound with $(b-a)(c-\tau)$. Taking the infimum over $s>0$ and using \cref{lem:testcurve-closed-fibers} proves \eqref{eq:finite-replacement-dual-bound}.

Since $v_t\leq ct$, taking $t\downarrow0$ in the bound
$\chi_\tau\leq v_t-t\tau$ gives $\chi_\tau\leq0$. The bounds in \eqref{eq:finite-replacement-dual-bound} show $\chi_\tau\sim\Gamma_\tau$. Hence it is a competitor for $P_\theta[\Gamma_\tau]=\Gamma_\tau$. Thus $\chi_\tau=\Gamma_\tau$ for every $\tau<c$. Recovering $v_t$ by \cref{lem:halfline-legendre-duality} proves $v_t=U_t$. This is maximality on every finite interval, including intervals starting at zero.
\end{proof}

\section{From a ray to a test curve}

In the reverse direction, the minimum principle produces the potentials of the test curve. To prove that they are envelope fixed points, we place a rooftop below a translated ray and use Perron maximality on intervals of arbitrary length. Shrinking the relative time then removes the translation error.

\begin{theorem}\label{thm:envelope-ray-to-testcurve}
If $u_t$ is an envelope ray of slope $c$, its inverse transform is an envelope test curve with cutoff $c$ and initial potential $\varphi$. The forward and inverse transforms are mutually inverse bijections between the two classes in \cref{def:hermitian-envelope-testcurve,def:hermitian-envelope-ray}.
\end{theorem}
\begin{proof}
By \cref{lem:invariant-strip-convexity}, each vertical function $u_t(x)-ct$ is finite and convex, or identically $-\infty$. A finite convex function bounded above on $(0,\infty)$ is decreasing, so $u_t-ct$ decreases everywhere. Since $\sup_X(u_t-ct)=0$, \cref{thm:normalized-compactness,cor:decreasing-limit} give a genuine decreasing limit as $t\to\infty$. Consequently
\[
 \widehat u_c=\inf_{t>0}(u_t-ct)\in\PSH(X,\theta).
\]
For $\tau<c$, the inverse transform is at least $\widehat u_c$ and at most $u_1-\tau$. Thus \cref{prop:hermitian-kiselman} makes it a genuine potential. Also
\[
 \widehat u_\tau\leq\inf_{t>0}t(c-\tau)=0
 \quad(\tau\leq c),
\]
while the same upper bound with $t\to\infty$ shows $\widehat u_\tau\equiv-\infty$ for $\tau>c$. Concavity and decreasingness in $\tau$ follow by taking an infimum of affine functions with slopes $-t$.

Write $g=\widehat u_\tau$ for a fixed $\tau<c$. We prove $P_\theta[g]=g$. For every $M>0$, the inequality $g\leq u_M-M\tau$ implies
\[
 P_\theta[g]\leq P_\theta[u_M]=\varphi.
\]
For $C>0$ put $q_C=P_\theta(g+C,0)$. It is genuine since $g\leq0$ is a competitor, and
\[
 g\leq q_C\leq P_\theta[g]\leq\varphi,
 \qquad q_C-C\leq g\leq u_M-M\tau.
\]
The affine path $q_C-Ct/M$ on $[0,M]$ is consequently admissible for $\varphi,u_M-M\tau$. By \cref{prop:hermitian-subgeodesic-calculus}, adding an affine function of time gives a bijection of Perron competitor families, so $u_t-t\tau$ is the Perron segment for these endpoints. For $0<s\leq1$ we obtain
\[
 q_C-Cs\leq u_{sM}-sM\tau.
\]
Take the infimum in $M>0$ to get $q_C-Cs\leq g$, and then let $s\downarrow0$. Thus $q_C=g$ for every $C>0$. Taking its increasing envelope limit proves $P_\theta[g]=g$.

Let $\psi=(\sup_{\tau<c}\widehat u_\tau)^*$. We have already shown $\widehat u_\tau\leq\varphi$, so $\psi\leq\varphi$. For almost every $x$, the convex function $t\mapsto u_t(x)$ has limit $\varphi(x)$ at zero. Indeed that limit exists by convexity and the linear upper bound, and an almost-everywhere convergent sequence from the $L^1$ endpoint condition identifies it. The endpoint identity in \cref{lem:halfline-legendre-duality} gives
\[
 \sup_{\tau<c}\inf_{t>0}(u_t(x)-t\tau)=\varphi(x).
\]
Therefore $\psi=\varphi$ by \cref{lem:qpsh-almost-everywhere-order}.

Finally, as $\tau\uparrow c$ the transforms decrease to $\widehat u_c$, because the two infima in $t$ and $\tau<c$ commute. Thus the endpoint convention agrees with \cref{def:hermitian-envelope-testcurve}. \Cref{lem:halfline-legendre-duality} recovers $u_t$ from its inverse transforms, including at points with identically $-\infty$ vertical functions. Conversely, \cref{lem:testcurve-closed-fibers} recovers a test curve from its forward transform. This proves the bijection.
\end{proof}

\Needspace{22\baselineskip}
\section{A logarithmic pole on the projective line}

On $\mathbb P^1$, a logarithmic pole gives a test curve that can be computed from tangent lines to a convex function. Its transform is a smooth ray, and the values at the pole show why the full parameter range is needed. In the standard affine coordinate $z$, take
\[
 \theta=\ddc\log(1+|z|^2),\qquad
 g(z)=\log\frac{|z|^2}{1+|z|^2}.
\]
The form extends to $\mathbb P^1$ and has integral one. The potential $g$ has a logarithmic pole at zero and extends smoothly across infinity, with value zero there.

Put $H(\lambda)=-\lambda\log\lambda-(1-\lambda)\log(1-\lambda)$ for $0<\lambda<1$, with $H(0)=H(1)=0$. For $0<\lambda<1$, define
\begin{equation}\label{eq:pone-explicit-model}
 \Gamma_\lambda(z)=
 \begin{cases}
 \lambda\log|z|^2-\log(1+|z|^2)+H(\lambda),
       & |z|^2\leq\lambda/(1-\lambda),\\
 0,    & |z|^2\geq\lambda/(1-\lambda).
 \end{cases}
\end{equation}
Set $\Gamma_\lambda=0$ for $\lambda\leq0$, $\Gamma_1=g$, and $\Gamma_\lambda=-\infty$ for $\lambda>1$.

\begin{proposition}\label{prop:pone-explicit-testcurve}
This is an envelope test curve with initial potential zero and cutoff one. Its forward transform is
\begin{equation}\label{eq:pone-explicit-ray}
 U_t(z)=\log\frac{1+\mathrm{e}^t|z|^2}{1+|z|^2},\qquad t>0.
\end{equation}
At $z=0$ its value is zero, and at infinity its value is $t$.
\end{proposition}
\begin{proof}
Fix $0<\lambda<1$. For $s=\log|z|^2$ let $q_\lambda(s)=\Gamma_\lambda(z)+\log(1+\mathrm{e}^s)$. In the first region of \eqref{eq:pone-explicit-model} it is the affine function $\lambda s+H(\lambda)$, tangent to $\log(1+\mathrm{e}^s)$ at $s_\lambda=\log(\lambda/(1-\lambda))$. In the second region it is $\log(1+\mathrm{e}^s)$. Its derivative is continuous and nondecreasing at the joining point, so $q_\lambda$ is convex. By \cref{cor:radial-subharmonic-convexity}, $\Gamma_\lambda$ is $\theta$-plurisubharmonic away from zero. It extends across its logarithmic pole by \cref{lem:qpsh-removable-analytic}. It is zero near infinity and has ordinary type $\lambda g$.

We show that it equals $P_\theta[\lambda g]$. For any competitor $v\leq0$ with $v\preceq\lambda g$, take the regularized supremum of its rotations. This remains a competitor, is radial, and dominates $v$. The bound $v\leq\lambda g+C$ is uniform under rotations, so the singularity condition is preserved by regularization. Write the resulting radial potential as $q(s)-\log(1+\mathrm{e}^s)$. By \cref{cor:radial-subharmonic-convexity}, $q$ is convex,
\[
 q(s)\leq\log(1+\mathrm{e}^s),\qquad q(s)\leq\lambda s+C
 \quad(s\to-\infty).
\]
The last inequality forces every secant slope of $q$ to be at least $\lambda$. Otherwise convexity would bound all slopes further to the left by a number smaller than $\lambda$, contradicting the displayed upper bound as $s\to-\infty$. For $s\leq s_\lambda$ it follows that
\[
 q(s)\leq q(s_\lambda)-\lambda(s_\lambda-s)
 \leq\lambda s+H(\lambda).
\]
For $s\geq s_\lambda$ the obstacle itself gives the required bound. Therefore $v\leq\Gamma_\lambda$. Since $\Gamma_\lambda$ is itself a competitor, it equals $P_\theta[\lambda g]$. Ordinary-type invariance now gives $P_\theta[\Gamma_\lambda]=\Gamma_\lambda$. Zero is also an envelope fixed point.

For a fixed $0<|z|<\infty$, put $p=|z|^2/(1+|z|^2)$. The function $\lambda\mapsto\Gamma_\lambda(z)$ is zero for $\lambda\leq p$. For $\lambda\geq p$ its derivative is
\[
 \log|z|^2+\log\frac{1-\lambda}{\lambda}.
\]
This starts at zero and decreases, proving concavity and decreasingness. At zero the curve is zero for $\lambda\leq0$ and $-\infty$ thereafter; at infinity it is zero for $\lambda\leq1$. These also satisfy the asserted conditions. Its limit towards $-\infty$ is zero, and its limit at one is $g$, proving that it is a test curve.

\Needspace{7\baselineskip}
For finite nonzero $z$ and $t>0$, the maximum in the transform occurs in the second parameter region at
\[
 \lambda=\frac{\mathrm{e}^t|z|^2}{1+\mathrm{e}^t|z|^2}.
\]
Substitution, or the identity for $a\in\mathbb R$
$\sup_{0\leq\lambda\leq1}(\lambda a+H(\lambda))=\log(1+\mathrm{e}^a)$,
gives \eqref{eq:pone-explicit-ray}. At zero only parameters $\lambda\leq0$ have finite values, giving zero; at infinity the supremum of $t\lambda$ for $\lambda<1$ is $t$. Thus the same formula extends smoothly in both charts.
\end{proof}

Each finite slice of this ray is bounded and has envelope zero, while the inverse transforms have logarithmic singularities. At $z=0$, the parameters $\lambda\leq0$ determine the transform: omitting them would incorrectly give $U_t(0)=-\infty$. Thus the parameter domain matters even when every positive parameter gives a pole.

\begin{figure}[htbp]
\centering
\resizebox{0.94\textwidth}{!}{\begin{tikzpicture}[font=\small,line width=0.7pt]
\begin{scope}[x=4.0cm,y=1.35cm]
\draw[->] (-0.08,0)--(1.1,0) node[right] {$\lambda$};
\draw[->] (0,-2.35)--(0,0.24) node[above] {$\Gamma_\lambda(z)$};
\foreach \x in {0.5,1} {\draw (\x,0.03)--(\x,-0.03) node[above,yshift=5pt] {$\x$};}
\foreach \y in {-1,-2} {\draw (-0.015,\y)--(0.015,\y) node[left,xshift=-5pt] {$\y$};}
\draw[solid] plot coordinates {(0.00000,0.00000) (0.00667,0.00000) (0.01333,0.00000) (0.02000,0.00000) (0.02667,0.00000) (0.03333,0.00000) (0.04000,0.00000) (0.04667,0.00000) (0.05333,0.00000) (0.06000,0.00000) (0.06667,0.00000) (0.07333,0.00000) (0.08000,0.00000) (0.08667,0.00000) (0.09333,0.00000) (0.10000,0.00000) (0.10667,0.00000) (0.11333,0.00000) (0.12000,-0.00000) (0.12667,-0.00026) (0.13333,-0.00092) (0.14000,-0.00196) (0.14667,-0.00338) (0.15333,-0.00515) (0.16000,-0.00726) (0.16667,-0.00970) (0.17333,-0.01246) (0.18000,-0.01553) (0.18667,-0.01891) (0.19333,-0.02257) (0.20000,-0.02653) (0.20667,-0.03075) (0.21333,-0.03525) (0.22000,-0.04002) (0.22667,-0.04504) (0.23333,-0.05032) (0.24000,-0.05585) (0.24667,-0.06162) (0.25333,-0.06763) (0.26000,-0.07387) (0.26667,-0.08035) (0.27333,-0.08705) (0.28000,-0.09397) (0.28667,-0.10112) (0.29333,-0.10849) (0.30000,-0.11606) (0.30667,-0.12385) (0.31333,-0.13185) (0.32000,-0.14006) (0.32667,-0.14847) (0.33333,-0.15708) (0.34000,-0.16589) (0.34667,-0.17490) (0.35333,-0.18411) (0.36000,-0.19351) (0.36667,-0.20310) (0.37333,-0.21289) (0.38000,-0.22286) (0.38667,-0.23303) (0.39333,-0.24338) (0.40000,-0.25392) (0.40667,-0.26464) (0.41333,-0.27555) (0.42000,-0.28664) (0.42667,-0.29791) (0.43333,-0.30936) (0.44000,-0.32100) (0.44667,-0.33281) (0.45333,-0.34481) (0.46000,-0.35698) (0.46667,-0.36934) (0.47333,-0.38187) (0.48000,-0.39458) (0.48667,-0.40747) (0.49333,-0.42054) (0.50000,-0.43378) (0.50667,-0.44720) (0.51333,-0.46080) (0.52000,-0.47458) (0.52667,-0.48854) (0.53333,-0.50267) (0.54000,-0.51698) (0.54667,-0.53148) (0.55333,-0.54615) (0.56000,-0.56100) (0.56667,-0.57603) (0.57333,-0.59124) (0.58000,-0.60664) (0.58667,-0.62221) (0.59333,-0.63797) (0.60000,-0.65392) (0.60667,-0.67005) (0.61333,-0.68636) (0.62000,-0.70286) (0.62667,-0.71956) (0.63333,-0.73644) (0.64000,-0.75351) (0.64667,-0.77078) (0.65333,-0.78824) (0.66000,-0.80589) (0.66667,-0.82375) (0.67333,-0.84180) (0.68000,-0.86006) (0.68667,-0.87852) (0.69333,-0.89719) (0.70000,-0.91606) (0.70667,-0.93515) (0.71333,-0.95445) (0.72000,-0.97397) (0.72667,-0.99372) (0.73333,-1.01368) (0.74000,-1.03387) (0.74667,-1.05429) (0.75333,-1.07495) (0.76000,-1.09585) (0.76667,-1.11699) (0.77333,-1.13838) (0.78000,-1.16002) (0.78667,-1.18192) (0.79333,-1.20409) (0.80000,-1.22653) (0.80667,-1.24924) (0.81333,-1.27224) (0.82000,-1.29553) (0.82667,-1.31913) (0.83333,-1.34303) (0.84000,-1.36726) (0.84667,-1.39181) (0.85333,-1.41671) (0.86000,-1.44196) (0.86667,-1.46759) (0.87333,-1.49359) (0.88000,-1.52000) (0.88667,-1.54683) (0.89333,-1.57411) (0.90000,-1.60185) (0.90667,-1.63008) (0.91333,-1.65884) (0.92000,-1.68816) (0.92667,-1.71808) (0.93333,-1.74866) (0.94000,-1.77996) (0.94667,-1.81205) (0.95333,-1.84501) (0.96000,-1.87898) (0.96667,-1.91412) (0.97333,-1.95064) (0.98000,-1.98889) (0.98667,-2.02945) (0.99333,-2.07355) (1.00000,-2.12693)};
\draw[dashed] plot coordinates {(0.00000,0.00000) (0.00667,0.00000) (0.01333,0.00000) (0.02000,0.00000) (0.02667,0.00000) (0.03333,0.00000) (0.04000,0.00000) (0.04667,0.00000) (0.05333,0.00000) (0.06000,0.00000) (0.06667,0.00000) (0.07333,0.00000) (0.08000,0.00000) (0.08667,0.00000) (0.09333,0.00000) (0.10000,0.00000) (0.10667,0.00000) (0.11333,0.00000) (0.12000,0.00000) (0.12667,0.00000) (0.13333,0.00000) (0.14000,0.00000) (0.14667,0.00000) (0.15333,0.00000) (0.16000,0.00000) (0.16667,0.00000) (0.17333,0.00000) (0.18000,0.00000) (0.18667,0.00000) (0.19333,0.00000) (0.20000,0.00000) (0.20667,0.00000) (0.21333,0.00000) (0.22000,0.00000) (0.22667,0.00000) (0.23333,0.00000) (0.24000,0.00000) (0.24667,0.00000) (0.25333,0.00000) (0.26000,0.00000) (0.26667,0.00000) (0.27333,0.00000) (0.28000,0.00000) (0.28667,0.00000) (0.29333,0.00000) (0.30000,0.00000) (0.30667,0.00000) (0.31333,0.00000) (0.32000,0.00000) (0.32667,0.00000) (0.33333,0.00000) (0.34000,0.00000) (0.34667,0.00000) (0.35333,0.00000) (0.36000,0.00000) (0.36667,0.00000) (0.37333,0.00000) (0.38000,0.00000) (0.38667,0.00000) (0.39333,0.00000) (0.40000,0.00000) (0.40667,0.00000) (0.41333,0.00000) (0.42000,0.00000) (0.42667,0.00000) (0.43333,0.00000) (0.44000,0.00000) (0.44667,0.00000) (0.45333,0.00000) (0.46000,0.00000) (0.46667,0.00000) (0.47333,0.00000) (0.48000,0.00000) (0.48667,0.00000) (0.49333,0.00000) (0.50000,0.00000) (0.50667,-0.00009) (0.51333,-0.00036) (0.52000,-0.00080) (0.52667,-0.00142) (0.53333,-0.00222) (0.54000,-0.00320) (0.54667,-0.00436) (0.55333,-0.00570) (0.56000,-0.00722) (0.56667,-0.00892) (0.57333,-0.01079) (0.58000,-0.01286) (0.58667,-0.01510) (0.59333,-0.01752) (0.60000,-0.02014) (0.60667,-0.02293) (0.61333,-0.02591) (0.62000,-0.02908) (0.62667,-0.03244) (0.63333,-0.03599) (0.64000,-0.03973) (0.64667,-0.04366) (0.65333,-0.04779) (0.66000,-0.05211) (0.66667,-0.05663) (0.67333,-0.06135) (0.68000,-0.06628) (0.68667,-0.07141) (0.69333,-0.07674) (0.70000,-0.08228) (0.70667,-0.08804) (0.71333,-0.09401) (0.72000,-0.10019) (0.72667,-0.10660) (0.73333,-0.11323) (0.74000,-0.12009) (0.74667,-0.12718) (0.75333,-0.13450) (0.76000,-0.14207) (0.76667,-0.14987) (0.77333,-0.15793) (0.78000,-0.16624) (0.78667,-0.17481) (0.79333,-0.18364) (0.80000,-0.19274) (0.80667,-0.20213) (0.81333,-0.21179) (0.82000,-0.22175) (0.82667,-0.23201) (0.83333,-0.24259) (0.84000,-0.25348) (0.84667,-0.26470) (0.85333,-0.27626) (0.86000,-0.28818) (0.86667,-0.30047) (0.87333,-0.31315) (0.88000,-0.32622) (0.88667,-0.33972) (0.89333,-0.35366) (0.90000,-0.36806) (0.90667,-0.38296) (0.91333,-0.39839) (0.92000,-0.41438) (0.92667,-0.43097) (0.93333,-0.44822) (0.94000,-0.46618) (0.94667,-0.48493) (0.95333,-0.50457) (0.96000,-0.52520) (0.96667,-0.54700) (0.97333,-0.57019) (0.98000,-0.59511) (0.98667,-0.62234) (0.99333,-0.65310) (1.00000,-0.69315)};
\draw[densely dotted] plot coordinates {(0.00000,0.00000) (0.00667,0.00000) (0.01333,0.00000) (0.02000,0.00000) (0.02667,0.00000) (0.03333,0.00000) (0.04000,0.00000) (0.04667,0.00000) (0.05333,0.00000) (0.06000,0.00000) (0.06667,0.00000) (0.07333,0.00000) (0.08000,0.00000) (0.08667,0.00000) (0.09333,0.00000) (0.10000,0.00000) (0.10667,0.00000) (0.11333,0.00000) (0.12000,0.00000) (0.12667,0.00000) (0.13333,0.00000) (0.14000,0.00000) (0.14667,0.00000) (0.15333,0.00000) (0.16000,0.00000) (0.16667,0.00000) (0.17333,0.00000) (0.18000,0.00000) (0.18667,0.00000) (0.19333,0.00000) (0.20000,0.00000) (0.20667,0.00000) (0.21333,0.00000) (0.22000,0.00000) (0.22667,0.00000) (0.23333,0.00000) (0.24000,0.00000) (0.24667,0.00000) (0.25333,0.00000) (0.26000,0.00000) (0.26667,0.00000) (0.27333,0.00000) (0.28000,0.00000) (0.28667,0.00000) (0.29333,0.00000) (0.30000,0.00000) (0.30667,0.00000) (0.31333,0.00000) (0.32000,0.00000) (0.32667,0.00000) (0.33333,0.00000) (0.34000,0.00000) (0.34667,0.00000) (0.35333,0.00000) (0.36000,0.00000) (0.36667,0.00000) (0.37333,0.00000) (0.38000,0.00000) (0.38667,0.00000) (0.39333,0.00000) (0.40000,0.00000) (0.40667,0.00000) (0.41333,0.00000) (0.42000,0.00000) (0.42667,0.00000) (0.43333,0.00000) (0.44000,0.00000) (0.44667,0.00000) (0.45333,0.00000) (0.46000,0.00000) (0.46667,0.00000) (0.47333,0.00000) (0.48000,0.00000) (0.48667,0.00000) (0.49333,0.00000) (0.50000,0.00000) (0.50667,0.00000) (0.51333,0.00000) (0.52000,0.00000) (0.52667,0.00000) (0.53333,0.00000) (0.54000,0.00000) (0.54667,0.00000) (0.55333,0.00000) (0.56000,0.00000) (0.56667,0.00000) (0.57333,0.00000) (0.58000,0.00000) (0.58667,0.00000) (0.59333,0.00000) (0.60000,0.00000) (0.60667,0.00000) (0.61333,0.00000) (0.62000,0.00000) (0.62667,0.00000) (0.63333,0.00000) (0.64000,0.00000) (0.64667,0.00000) (0.65333,0.00000) (0.66000,0.00000) (0.66667,0.00000) (0.67333,0.00000) (0.68000,0.00000) (0.68667,0.00000) (0.69333,0.00000) (0.70000,0.00000) (0.70667,0.00000) (0.71333,0.00000) (0.72000,0.00000) (0.72667,0.00000) (0.73333,0.00000) (0.74000,0.00000) (0.74667,0.00000) (0.75333,0.00000) (0.76000,0.00000) (0.76667,0.00000) (0.77333,0.00000) (0.78000,0.00000) (0.78667,0.00000) (0.79333,0.00000) (0.80000,0.00000) (0.80667,0.00000) (0.81333,0.00000) (0.82000,0.00000) (0.82667,0.00000) (0.83333,0.00000) (0.84000,0.00000) (0.84667,0.00000) (0.85333,0.00000) (0.86000,0.00000) (0.86667,0.00000) (0.87333,0.00000) (0.88000,0.00000) (0.88667,-0.00017) (0.89333,-0.00077) (0.90000,-0.00185) (0.90667,-0.00341) (0.91333,-0.00550) (0.92000,-0.00816) (0.92667,-0.01142) (0.93333,-0.01533) (0.94000,-0.01996) (0.94667,-0.02538) (0.95333,-0.03168) (0.96000,-0.03898) (0.96667,-0.04745) (0.97333,-0.05730) (0.98000,-0.06889) (0.98667,-0.08278) (0.99333,-0.10021) (1.00000,-0.12693)};
\end{scope}
\begin{scope}[xshift=6.1cm,x=1.0cm,y=0.76cm]
\draw[->] (0,0)--(4.4,0) node[right] {$t$};
\draw[->] (0,0)--(0,4.4) node[above] {$U_t(z)$};
\foreach \x in {2,4} {\draw (\x,0.05)--(\x,-0.05) node[below] {$\x$};}
\foreach \y in {2,4} {\draw (-0.04,\y)--(0.04,\y) node[left] {$\y$};}
\draw[solid] plot coordinates {(0.00000,0.00000) (0.02667,0.00322) (0.05333,0.00651) (0.08000,0.00988) (0.10667,0.01333) (0.13333,0.01686) (0.16000,0.02047) (0.18667,0.02417) (0.21333,0.02795) (0.24000,0.03182) (0.26667,0.03578) (0.29333,0.03983) (0.32000,0.04397) (0.34667,0.04821) (0.37333,0.05254) (0.40000,0.05697) (0.42667,0.06150) (0.45333,0.06613) (0.48000,0.07087) (0.50667,0.07570) (0.53333,0.08065) (0.56000,0.08570) (0.58667,0.09087) (0.61333,0.09614) (0.64000,0.10153) (0.66667,0.10703) (0.69333,0.11266) (0.72000,0.11840) (0.74667,0.12426) (0.77333,0.13024) (0.80000,0.13635) (0.82667,0.14259) (0.85333,0.14896) (0.88000,0.15545) (0.90667,0.16208) (0.93333,0.16884) (0.96000,0.17573) (0.98667,0.18277) (1.01333,0.18994) (1.04000,0.19725) (1.06667,0.20470) (1.09333,0.21230) (1.12000,0.22005) (1.14667,0.22794) (1.17333,0.23598) (1.20000,0.24417) (1.22667,0.25252) (1.25333,0.26101) (1.28000,0.26967) (1.30667,0.27848) (1.33333,0.28744) (1.36000,0.29657) (1.38667,0.30586) (1.41333,0.31530) (1.44000,0.32492) (1.46667,0.33469) (1.49333,0.34464) (1.52000,0.35475) (1.54667,0.36502) (1.57333,0.37547) (1.60000,0.38609) (1.62667,0.39687) (1.65333,0.40783) (1.68000,0.41896) (1.70667,0.43027) (1.73333,0.44175) (1.76000,0.45340) (1.78667,0.46523) (1.81333,0.47724) (1.84000,0.48942) (1.86667,0.50177) (1.89333,0.51431) (1.92000,0.52702) (1.94667,0.53991) (1.97333,0.55297) (2.00000,0.56622) (2.02667,0.57964) (2.05333,0.59324) (2.08000,0.60702) (2.10667,0.62097) (2.13333,0.63511) (2.16000,0.64942) (2.18667,0.66390) (2.21333,0.67856) (2.24000,0.69340) (2.26667,0.70842) (2.29333,0.72360) (2.32000,0.73896) (2.34667,0.75450) (2.37333,0.77021) (2.40000,0.78609) (2.42667,0.80214) (2.45333,0.81836) (2.48000,0.83475) (2.50667,0.85130) (2.53333,0.86803) (2.56000,0.88492) (2.58667,0.90197) (2.61333,0.91919) (2.64000,0.93657) (2.66667,0.95411) (2.69333,0.97181) (2.72000,0.98967) (2.74667,1.00768) (2.77333,1.02585) (2.80000,1.04417) (2.82667,1.06265) (2.85333,1.08127) (2.88000,1.10005) (2.90667,1.11897) (2.93333,1.13804) (2.96000,1.15725) (2.98667,1.17660) (3.01333,1.19610) (3.04000,1.21573) (3.06667,1.23550) (3.09333,1.25541) (3.12000,1.27545) (3.14667,1.29562) (3.17333,1.31592) (3.20000,1.33635) (3.22667,1.35691) (3.25333,1.37759) (3.28000,1.39840) (3.30667,1.41932) (3.33333,1.44037) (3.36000,1.46153) (3.38667,1.48281) (3.41333,1.50420) (3.44000,1.52570) (3.46667,1.54732) (3.49333,1.56904) (3.52000,1.59087) (3.54667,1.61280) (3.57333,1.63484) (3.60000,1.65697) (3.62667,1.67921) (3.65333,1.70154) (3.68000,1.72397) (3.70667,1.74650) (3.73333,1.76911) (3.76000,1.79182) (3.78667,1.81462) (3.81333,1.83750) (3.84000,1.86047) (3.86667,1.88353) (3.89333,1.90666) (3.92000,1.92988) (3.94667,1.95318) (3.97333,1.97655) (4.00000,2.00000)};
\draw[dashed] plot coordinates {(0.00000,0.00000) (0.02667,0.01342) (0.05333,0.02702) (0.08000,0.04080) (0.10667,0.05475) (0.13333,0.06889) (0.16000,0.08320) (0.18667,0.09768) (0.21333,0.11234) (0.24000,0.12718) (0.26667,0.14220) (0.29333,0.15738) (0.32000,0.17275) (0.34667,0.18828) (0.37333,0.20399) (0.40000,0.21987) (0.42667,0.23592) (0.45333,0.25214) (0.48000,0.26853) (0.50667,0.28508) (0.53333,0.30181) (0.56000,0.31870) (0.58667,0.33575) (0.61333,0.35297) (0.64000,0.37035) (0.66667,0.38789) (0.69333,0.40559) (0.72000,0.42345) (0.74667,0.44146) (0.77333,0.45963) (0.80000,0.47795) (0.82667,0.49643) (0.85333,0.51505) (0.88000,0.53383) (0.90667,0.55275) (0.93333,0.57182) (0.96000,0.59103) (0.98667,0.61038) (1.01333,0.62988) (1.04000,0.64951) (1.06667,0.66928) (1.09333,0.68919) (1.12000,0.70923) (1.14667,0.72940) (1.17333,0.74970) (1.20000,0.77014) (1.22667,0.79069) (1.25333,0.81137) (1.28000,0.83218) (1.30667,0.85310) (1.33333,0.87415) (1.36000,0.89531) (1.38667,0.91659) (1.41333,0.93798) (1.44000,0.95948) (1.46667,0.98110) (1.49333,1.00282) (1.52000,1.02465) (1.54667,1.04658) (1.57333,1.06862) (1.60000,1.09075) (1.62667,1.11299) (1.65333,1.13532) (1.68000,1.15775) (1.70667,1.18028) (1.73333,1.20290) (1.76000,1.22560) (1.78667,1.24840) (1.81333,1.27128) (1.84000,1.29425) (1.86667,1.31731) (1.89333,1.34044) (1.92000,1.36366) (1.94667,1.38696) (1.97333,1.41033) (2.00000,1.43378) (2.02667,1.45731) (2.05333,1.48090) (2.08000,1.50457) (2.10667,1.52831) (2.13333,1.55212) (2.16000,1.57600) (2.18667,1.59994) (2.21333,1.62395) (2.24000,1.64802) (2.26667,1.67215) (2.29333,1.69634) (2.32000,1.72059) (2.34667,1.74490) (2.37333,1.76927) (2.40000,1.79369) (2.42667,1.81816) (2.45333,1.84269) (2.48000,1.86727) (2.50667,1.89191) (2.53333,1.91659) (2.56000,1.94132) (2.58667,1.96609) (2.61333,1.99091) (2.64000,2.01578) (2.66667,2.04070) (2.69333,2.06565) (2.72000,2.09065) (2.74667,2.11569) (2.77333,2.14077) (2.80000,2.16589) (2.82667,2.19104) (2.85333,2.21624) (2.88000,2.24147) (2.90667,2.26674) (2.93333,2.29204) (2.96000,2.31737) (2.98667,2.34274) (3.01333,2.36815) (3.04000,2.39358) (3.06667,2.41904) (3.09333,2.44454) (3.12000,2.47006) (3.14667,2.49562) (3.17333,2.52120) (3.20000,2.54681) (3.22667,2.57244) (3.25333,2.59810) (3.28000,2.62379) (3.30667,2.64950) (3.33333,2.67524) (3.36000,2.70100) (3.38667,2.72678) (3.41333,2.75259) (3.44000,2.77841) (3.46667,2.80426) (3.49333,2.83013) (3.52000,2.85602) (3.54667,2.88193) (3.57333,2.90786) (3.60000,2.93381) (3.62667,2.95978) (3.65333,2.98576) (3.68000,3.01176) (3.70667,3.03778) (3.73333,3.06382) (3.76000,3.08987) (3.78667,3.11594) (3.81333,3.14202) (3.84000,3.16812) (3.86667,3.19423) (3.89333,3.22036) (3.92000,3.24650) (3.94667,3.27265) (3.97333,3.29882) (4.00000,3.32500)};
\draw[densely dotted] plot coordinates {(0.00000,0.00000) (0.02667,0.02353) (0.05333,0.04712) (0.08000,0.07079) (0.10667,0.09453) (0.13333,0.11834) (0.16000,0.14222) (0.18667,0.16616) (0.21333,0.19017) (0.24000,0.21424) (0.26667,0.23837) (0.29333,0.26256) (0.32000,0.28681) (0.34667,0.31112) (0.37333,0.33549) (0.40000,0.35991) (0.42667,0.38438) (0.45333,0.40891) (0.48000,0.43349) (0.50667,0.45812) (0.53333,0.48281) (0.56000,0.50753) (0.58667,0.53231) (0.61333,0.55713) (0.64000,0.58200) (0.66667,0.60691) (0.69333,0.63187) (0.72000,0.65687) (0.74667,0.68191) (0.77333,0.70699) (0.80000,0.73210) (0.82667,0.75726) (0.85333,0.78246) (0.88000,0.80769) (0.90667,0.83295) (0.93333,0.85826) (0.96000,0.88359) (0.98667,0.90896) (1.01333,0.93436) (1.04000,0.95980) (1.06667,0.98526) (1.09333,1.01076) (1.12000,1.03628) (1.14667,1.06184) (1.17333,1.08742) (1.20000,1.11303) (1.22667,1.13866) (1.25333,1.16432) (1.28000,1.19001) (1.30667,1.21572) (1.33333,1.24146) (1.36000,1.26722) (1.38667,1.29300) (1.41333,1.31881) (1.44000,1.34463) (1.46667,1.37048) (1.49333,1.39635) (1.52000,1.42224) (1.54667,1.44815) (1.57333,1.47408) (1.60000,1.50003) (1.62667,1.52600) (1.65333,1.55198) (1.68000,1.57798) (1.70667,1.60400) (1.73333,1.63004) (1.76000,1.65609) (1.78667,1.68216) (1.81333,1.70824) (1.84000,1.73434) (1.86667,1.76045) (1.89333,1.78658) (1.92000,1.81272) (1.94667,1.83887) (1.97333,1.86504) (2.00000,1.89122) (2.02667,1.91742) (2.05333,1.94362) (2.08000,1.96984) (2.10667,1.99607) (2.13333,2.02231) (2.16000,2.04856) (2.18667,2.07482) (2.21333,2.10109) (2.24000,2.12738) (2.26667,2.15367) (2.29333,2.17997) (2.32000,2.20628) (2.34667,2.23261) (2.37333,2.25894) (2.40000,2.28527) (2.42667,2.31162) (2.45333,2.33798) (2.48000,2.36434) (2.50667,2.39071) (2.53333,2.41709) (2.56000,2.44348) (2.58667,2.46987) (2.61333,2.49628) (2.64000,2.52268) (2.66667,2.54910) (2.69333,2.57552) (2.72000,2.60195) (2.74667,2.62838) (2.77333,2.65482) (2.80000,2.68127) (2.82667,2.70772) (2.85333,2.73418) (2.88000,2.76064) (2.90667,2.78711) (2.93333,2.81358) (2.96000,2.84006) (2.98667,2.86654) (3.01333,2.89303) (3.04000,2.91952) (3.06667,2.94602) (3.09333,2.97252) (3.12000,2.99903) (3.14667,3.02554) (3.17333,3.05205) (3.20000,3.07857) (3.22667,3.10510) (3.25333,3.13162) (3.28000,3.15815) (3.30667,3.18468) (3.33333,3.21122) (3.36000,3.23776) (3.38667,3.26431) (3.41333,3.29085) (3.44000,3.31740) (3.46667,3.34396) (3.49333,3.37051) (3.52000,3.39707) (3.54667,3.42363) (3.57333,3.45020) (3.60000,3.47676) (3.62667,3.50333) (3.65333,3.52991) (3.68000,3.55648) (3.70667,3.58306) (3.73333,3.60964) (3.76000,3.63622) (3.78667,3.66280) (3.81333,3.68939) (3.84000,3.71598) (3.86667,3.74257) (3.89333,3.76916) (3.92000,3.79575) (3.94667,3.82235) (3.97333,3.84895) (4.00000,3.87555)};
\end{scope}
\end{tikzpicture}
}
\caption{The explicit test curve and its ray, at points with $\log|z|^2=-2$ (solid), $0$ (dashed), and $2$ (dotted). The parameter functions on the left are concave; their transforms on the right are convex.}
\label{fig:explicit-legendre-pone}
\end{figure}
